\section{Introduction}\label{sec:intro}

\paragraph{The question.} Let $b(d)=\max(1,\lceil\log_2d\rceil)$ and $L=d^2b(d)$. Statement~(T) asserts that one
deterministic multitape Turing machine transposes every $d\times d$ matrix with $b(d)$-bit entries, given and returned in
row-major order, within $O(L)$ steps. The obvious recursive algorithm takes $O(L\log d)$ steps. By a reduction of Harvey
and van der Hoeven, discussed in the first paper of the series~\cite{FirstPaper2026}, a linear-time multiplication
machine would give a linear-time transposer, so a proof of $\neg$(T) would give a superlinear lower bound for integer
multiplication on multitape Turing machines. Whether (T) holds is open, and this paper does not decide it. It is the
eleventh paper of a series that collects the necessary conditions a linear-time transposer must satisfy and the
constructions they leave open.

\paragraph{Where the series stands.} The fourth paper reduced the impossibility of oblivious linear-time transposers to
two statements~\cite{FourthPaper2026}: \emph{depth reduction}, that the dependency DAG of every linear-time schedule can
be cut to bounded depth by a cheap genie, and a coding bound on the remaining bounded-depth cones, the \emph{crux}. On the
depth side, the eighth paper's Theorem~F cuts the walk DAG of the work heads to bounded depth whenever the heads have
bounded owner-load, in particular whenever they enter every cell boundedly
often~\cite[Theorem~3.6 and Corollary~3.2]{BoundedNesting2026}; whether such cuts exist for all pairs of walks is the open statement
$\TCP(2)$, and the deep-nesting pairs are its sharpest test case~\cite{BoundedNesting2026,DirectedReciprocity2026}. On the
coding side, the fifth paper showed that the cone relations of every linear-time schedule are paired-distance separated
after a paid trimming, and its Question~9.1 asks whether every storage code on such relations of bounded degree has rate
at most one half~\cite{DirectedRecovery2026}. The answer is affirmative on several classes of
relations~\cite{DirectedRecovery2026,SeventhPaper2026,DirectedReciprocity2026}; we write [Q9.1] for an affirmative answer in
general. Together with Theorem~F, [Q9.1] excludes oblivious transposers whose work heads enter every cell boundedly often
(the eighth paper's Corollary~F$'$~\cite[Corollary~6.3]{BoundedNesting2026}). Independently of coding hypotheses, the
ninth paper excludes oblivious geometries with bounded span overlap and few large return families~\cite{SourceCovers2026},
and the frozen gap-area certificate cannot exclude the remaining two-exception inverse-square case at moduli comparable to
the protected rank~\cite{GapAreaMemo2026}. For adaptive machines, the eighth paper's Corollary~F$''$ extends Corollary~F$'$,
given [Q9.1], to bounded-entry machines whose traces carry at most $\eta L$ bits for a fixed
$\eta<\frac12$~\cite[Corollary~6.4]{BoundedNesting2026}; padding the motion of the input head defeats that
hypothesis without changing the computation~\cite[Section~6.2]{BoundedNesting2026}.

\paragraph{This paper.} Given [Q9.1], we show that every linear-time transposer revisits some tape unboundedly often,
that an oblivious one must in addition merge deep lineages unboundedly often in its code, and that one whose work tapes
have bounded entry must steer its input head by the data. Unconditionally, at block scale every linear-time transposer,
adaptive or not, runs an oblivious external-memory transposition with linearly many I/Os on a large set of inputs, so
that one coding problem underlies all cases.
\begin{enumerate}
\item \emph{Merge depth replaces schedule depth} (\cref{sec:hop}). The eighth paper's Corollary~F cuts the whole schedule
to bounded depth with Theorem~F, which needs bounded owner-load; schedules that revisit cells unboundedly often, such as
the deep-nesting pairs, may have no such cut. Hop charging handles them. Theorem~F cuts only the hops at boundedly entered
cells (the cold genie, \cref{lem:cold}); the hops at the other cells, the deep hops, are charged instead, one revealed
coordinate for every $\tau$ deep hops along a lineage of the code. If the code merges deep lineages at most $s$ times
along every path from most outputs, then, on \emph{every} schedule, however deeply nested and however many cells it
revisits, the transposition matrix splits into a part of small rank and a part supported on balls of a radius that does
not depend on the schedule (Theorems~HC and~HC$'$, for linear codes), and for all codes the side information inside
such balls has a correspondingly short index code (Theorem~HC$'$(b)). So the linear crux excludes such transposers with
linear codes (Corollary~HC), and [Q9.1] excludes them for all codes (Corollary~HC$^*$). This answers, under these
hypotheses and for codes of bounded merge depth, the eighth paper's question whether a scratch of revisited cells can
help~\cite[Problem~10.5]{BoundedNesting2026}. Routing through deep hops is limited by their number
(\cref{prop:hop-routing}), the number of deep cells adds nothing (\cref{prop:capacity}), and a cheap cut of the code's
merge structure would suffice (\cref{prop:merge-reduction}). What remains for oblivious machines is depth reduction of
that structure (Conjecture~HV$^*$, \cref{conj:HV}).
\item \emph{Skeleton entropy and input steering} (\cref{sec:skeleton}). Theorem~SE is the form of Corollary~F$''$ that
is robust against padding of bounded entry, as the eighth paper asks: it bounds the entropy of the coarse skeleton at a
fixed scale~$w$, the head positions and the output count rounded down to multiples of~$w$, instead of the entropy of the
trace, and it contains Corollary~F$''$. Its proof turns the machine into a \emph{block machine} that is oblivious on
each skeleton class, so that Theorem~F, the cone lemma and the paired-distance argument of the fifth paper apply class
by class. Heads of bounded entry carry almost no coarse entropy at large scales, and two corollaries follow, given
[Q9.1]. No linear-time transposer, adaptive or not, has bounded entry on every tape, the input tape included
(Corollary~RV). A transposer whose work tapes have bounded entry must therefore revisit input cells unboundedly often,
and its input head must carry about $L/2$ bits of the data in its coarse motion at every fixed scale (Corollary~IS,
which contains Corollary~F$'$ and extends it to adaptive machines). Padding of bounded entry before the run does not
help a machine meet this condition, but padding that revisits input cells unboundedly often meets it at no cost
(\cref{rem:skel-padding}): it is a necessary condition, not an obstruction.
\item \emph{Localization} (\cref{sec:localization}). Theorem~L combines these. Given [Q9.1], every linear-time
transposer revisits some tape unboundedly often; if it is oblivious, it revisits a work tape and its code merges deep
lineages unboundedly often, for every pruning, cold genie and lineage assignment; and if its work tapes have bounded
entry, it revisits input cells and carries about half of the input in the coarse motion of its input head. Three kinds
of machines remain, and the open problems concern them.
\item \emph{Unconditional reductions} (\cref{sec:reductions}). A machine is excluded if, for a small enough $\eta>0$, it
has a trace of probability at least $2^{-\eta L}$ whose return families fall under the source-cover exclusions
(Theorem~LA); the passage from a machine to the ninth paper's admissible geometries is formal. At block scale every
linear-time transposer, adaptive or not, yields an external-memory transposition with $O(d)$ I/Os, linear in the number
of blocks, that is oblivious on a set of $2^{L-O(d)}$ inputs (Theorem~EM). External-memory lower bounds of this kind are
known only for algorithms that move records as indivisible atoms~\cite{AggarwalVitter1988}, or conditionally on a
network-coding conjecture~\cite{FHLS2019} that is false in general~\cite{BravermanHe2026}; an indivisibility-free bound
for this oblivious, large-set form of transposition would refute (T) for every number of tapes (\cref{cor:EM}).
\item \emph{Deep nesting} (\cref{sec:dnp}). The eighth paper's Conjecture~A$''$ would make the deep-nesting pairs
depth-robust and so refute $\TCP(2)$. Theorem~BP proves it for pairs of block labels of bounded length at one odd shift,
with no $\log q$ term and no cost hypothesis, conditional on Conjecture~S (\cref{conj:S}): the straddle counts of most
coarse classes concentrate as the shift varies. We give partial results on Conjecture~S: an exact Fourier form of the
normalized variances, a family of classes on which they stay large, and bounds that reduce the conjecture to the average
of their non-local part. Theorem~BP is evidence that the deep-nesting pairs are depth-robust; it does not decide
$\TCP(2)$, since general labels remain.
\item \emph{What a construction would need} (\cref{sec:construction}). By the first paper, every linear-time transposer
carries, on the block network of its run, a causal code that solves complete exchange on a large set of inputs and
beats multicommodity flow by a factor $\Omega(\log d)$~\cite[Theorems~8.8 and~9.2]{FirstPaper2026}. \Cref{sec:construction} collects what the series now requires of
such a machine and states the two forms that remain: an oblivious machine whose code merges deep lineages along patterns
without cheap cuts, or an adaptive machine whose revisits the present methods do not constrain further. Lemma~PG rules
out one natural building block: Tanner codes assembled from projective planes, with every coordinate recovered from one
line of a plane, violate the paired-distance condition for every order $q\ge7$.
\end{enumerate}

\paragraph{How to read this paper: the formal proofs.} Every numbered theorem, lemma, corollary and proposition is
formalized in Lean~4 with Mathlib, and every declaration that the paper cites depends at most on the axioms
\texttt{propext}, \texttt{Classical.choice} and \texttt{Quot.sound}. Each numbered result carries a status line that names
its declarations and their file, and numbered definitions carry one when they have a Lean counterpart. The proof of a
numbered result is an outline of the formal proof; the Lean statement is authoritative, and the outline says where it
is stronger than the printed statement, for instance in its parameter ranges or because it uses only part of the
hypotheses. Results that assume an affirmative answer to Question~9.1, the linear crux or \cref{conj:S} say so in their
statements, and their Lean versions take these as explicit hypotheses: \lean{MergeDepth.Q91}, or its binary form
\lean{MergeDepth.Q91Binary}, \lean{MergeDepth.LinCrux} and \lean{MergeDepth.ConjectureS}. The formal proofs build on the
Lean libraries of the eighth, fifth, ninth and tenth papers, in which the results they use, Theorem~F among them, are
Lean-verified; what they need from results that earlier papers formalize only in part or not at all, such as the cone
lemma, the paired-distance theorem and the uniformity argument of the fifth paper's Proposition~9.2, is reproved in Lean
in the forms used here. Results of earlier papers are cited by their numbers at a pinned revision. Claims that have only a
proof sketch appear in remarks marked ``Proof sketch only'', numerical observations in remarks marked as numerical, and
heuristics are labelled as such; none of them is used in a proof. \Cref{sec:leanindex} lists the declarations, and the
paper's README explains how to build the library and check its axioms.

\paragraph{Related work.} The machine model, the dependency DAGs, genies, the cone lemma and the crux are those of the
fourth and eighth papers~\cite{FourthPaper2026,BoundedNesting2026}; the paired-distance theorem and Question~9.1 are from
the fifth~\cite{DirectedRecovery2026}; the source-cover exclusions are from the ninth~\cite{SourceCovers2026}; the
deep-nesting results build on the tenth~\cite{DirectedReciprocity2026}. Theorem~EM is a standard simulation of tapes in
the external-memory model of Aggarwal and Vitter~\cite{AggarwalVitter1988}, followed by the popular-trace argument that
also underlies the non-oblivious case of~\cite{FHLS2019}; its point here is the conclusion for transposition with
internal memory of a constant number of blocks. The network-coding conjecture behind the conditional bounds
of~\cite{FHLS2019} is false: Braverman and He showed that for multiple unicast in undirected networks, coding can beat
multicommodity flow by a factor that grows with the size of the network~\cite{BravermanHe2026}.

\paragraph{Organization.} \Cref{sec:prelim} recalls the notation. \Cref{sec:hop,sec:skeleton} prove the hop-charging and
skeleton-entropy theorems, and \cref{sec:localization} combines them. \Cref{sec:reductions} proves the unconditional
reductions, and \cref{sec:dnp} the block-pair theorem and the partial results on Conjecture~S. \Cref{sec:construction}
collects what a construction would need and proves Lemma~PG, and \cref{sec:open} lists the open problems.
\Cref{sec:leanindex} indexes the Lean declarations.
